\section{Continued recovery, implementation and observational experiments}
\label{sec:continuation}

The earlier assessment identified missing mechanisms too broadly. Reading the
complete Tanaka thesis and the 2022 scholarly preprint recovered the DOECLIM
constants, the coastal decision menu and the socioeconomic priors. Independent
providers supplied forcing, temperature, carbon emissions, ocean heat and
coastal-hazard inputs \cite{inputs}. The active audit in
\texttt{docs/BLOCKER\_AUDIT.md} distinguishes recovered specifications,
derived corrections, explicit alternatives and remaining data gaps.

\subsection{Coupled energy balance and numerical heat transport}

The land--troposphere and ocean--troposphere boxes exchange heat and lose energy
through separate radiative feedbacks \cite{tanaka,urban}. For land fraction
$f=0.29$, their temperatures $T_L,T_S$ satisfy
\begin{align}
C_L\dot T_L&=Q-\lambda_LT_L-\frac{k}{f}(T_L-bT_S),\\
C_S\dot T_S&=Q-\lambda_ST_S+\frac{k}{1-f}(T_L-bT_S)-F_O.
\end{align}
Here $b=1.3$, $C_L=0.52$ and $C_S=7.8$\,W\,yr\,m$^{-2}$\,K$^{-1}$.
The global surface-air observable is $fT_L+(1-f)bT_S$. The equilibrium land/SST
ratio is 1.43 and $k=1.59-0.31\lambda_L$. Together with an equilibrium warming
of ECS under 3.7\,W\,m$^{-2}$, these relations determine both feedbacks and the
exchange coefficient. The explicit formulas are recorded in
\texttt{docs/RECOVERED\_SPECIFICATIONS.md}.

An 80-cell, 4000\,m interior ocean exchanges heat continuously with the mixed
layer and has an insulated bottom. Finite-volume fluxes cancel under a global
heat-capacity weighted sum; this is tested independently of equilibrium fitting.
Crank--Nicolson integration uses four steps per year. Refining the vertical grid
from 80 to 160 cells changes the year-70 warming under 1\% annual CO$_2$ growth
by 0.000787\,K; the next doubling changes it by 0.000197\,K. This approximately
fourfold error reduction agrees with second-order spatial convergence.

The 60\,m effective mixed-layer capacity contains about 2.5\,m of
atmospheric-equivalent capacity. Comparisons with NOAA integrate only 57.5\,m
of water and the appropriate fractional interior cells down to the stated
observational depth. Whole-column heat is not compared to a 700\,m record.
HadCRUT4 blends land air temperature with SST, whereas the legacy DOECLIM
observable includes enhanced marine-air warming. Treating that blended record
as the global-air observable follows the simplified legacy calibration and is
an additional representation approximation, not an exact observation operator.

\subsection{A closed, explicitly alternative carbon cycle}

The recovered terrestrial stocks are 100, 500, 120 and 1500\,Gt\,C in leaves,
wood, detritus and soil \cite{siegenthaler}. Baseline NPP of 60\,Gt\,C\,yr$^{-1}$
allocates 35 to leaves and 25 to wood. Leaf turnover sends 35 to detritus;
wood sends 20 to detritus and 5 to soil. Detritus respires 45 and transfers 10
to soil; soil respires 15. Linear turnover rates are these fluxes divided by
their source stocks. Fertilization modifies NPP and the documented $Q_{10}$
factor modifies the two respiration fluxes \cite{ricciuto}. All internal
transfers cancel in the global carbon balance.

Ocean carbon uses an instantaneously equilibrated atmosphere--mixed-layer
reservoir and three deeper reservoirs. Constant alkalinity, total boron and
carbonate equilibrium determine atmospheric CO$_2$ and mixed-layer carbon from
their combined inventory \cite{hooss1999}. The nonlinear scalar root is solved
by bisection in log hydrogen concentration. At 280\,ppm the reconstructed DIC
is 0.002080909\,mol\,L$^{-1}$, consistent with the printed H99 value to its
reported precision. The mixed-layer depth and deep exchange parameters are
from the later Tanaka description. Mixing these documented versions, and
holding carbonate chemistry temperature fixed, defines an alternative closure;
it does not establish exact recovery of the target SNEASY version.

Second-order splitting couples carbon RK4 steps to climate CN steps. Twenty-six
checks establish baseline equilibrium, chemical partition, stock positivity,
second-order coupled convergence and total-carbon conservation. Across 18
historical sensitivities, the largest carbon-budget residual is below
$10^{-11}$\,Gt\,C. Conservation alone does not validate the response: the
uncalibrated default predicts 461.38\,ppm in 2022, compared with NOAA's
417.08\,ppm. The GCB historical table has no pre-1850 land-use emissions; the
input retains blanks, while separate calculations prescribe zero or the 1850
rate before 1850. These are alternatives, not recovered observations.

\begin{figure}[htbp]
\centering\includegraphics[width=\textwidth]{carbon_alternative.pdf}
\caption{Uncalibrated carbon--climate experiments, not target reproductions.
Left: historical CO$_2$ sensitivities to fertilization and the explicitly
assumed early land-use input, with independent NOAA observations. Centre:
documented emissions recurrence initialized from GCB2023, with three prescribed
peak years. Right: coupled warming under fixed SSP2-4.5 non-CO$_2$ forcing.
The historical mismatch limits interpretation of these future curves.}
\end{figure}

\subsection{Temperature-only calibration and failed holdout checks}

The first observational experiment estimates ECS, heat diffusivity, aerosol
multiplier, temperature offset, AR1 innovation scale and AR1 correlation.
HadCRUT4.6.0.0 years 1850--2005 form the training data; years 2006--2021 are
held out. Known observational standard deviations are approximated by the total
95\% interval widths divided by $2(1.96)$. This approximation omits the record's
full temporally correlated measurement covariance. A scalar Kalman filter
integrates the latent AR1 residual exactly under these assumptions; tests match
the dense Gaussian covariance likelihood to $10^{-12}$.

Uniform bounds follow the B17 climate table, except that ECS is restricted to
0.7--10\,K: the recovered exchange relation becomes negative below about
0.68\,K. CEDA effective forcing supplies historical and future inputs, with
both aerosol components multiplied by the same parameter. This differs from
the older indirect-aerosol convention and from Darnell's uncertain emissions
and nearest-SSP forcing construction. Neither the posterior nor the six
scenario experiments is claimed as the Wong ensemble.

An initial short run failed convergence checks. The retained calculation uses
four dispersed chains, each with 6000 warmup and 30,000 production steps;
every fifth state is retained. Proposal adaptation stops before production.
Log ECS and log diffusivity proposals include their Jacobian, preserving the
physical-coordinate prior. Independently implemented rank/folded split
$\widehat R$ and bulk/tail ESS follow Vehtari et al.\ \cite{vehtari}.
The largest $\widehat R$ is 1.00427 and the smallest ESS is 1894, satisfying
the prespecified limits of 1.01 and 400.

Conditional forecasts propagate the final training-period residual state and
include both AR1 innovation and observational uncertainty. Despite sampler
convergence, holdout temperature RMSE is 0.209\,K, with mean bias
$+0.189$\,K. Only 10 of 16 held-out annual temperatures lie inside the nominal
90\% predictive intervals. After fitting only an additive heat datum offset
over 1955--2005, the 700\,m heat parameter-median RMSE in 2006--2021 is
$16.46\times10^{22}$\,J. These failures motivate a separate temperature--heat
calibration; they are not removed from the record.

One thousand complete joint parameter rows generate six SSP trajectories;
250- and 500-row subsets test sampling variation. The physical temperature
intervals exclude future residual noise. SSP2-4.5 mean warming over
2081--2100 is 2.924\,K, with a parameter 90\% interval of 1.976--4.188\,K.
The corresponding printed AR6 assessment is 2.7 [2.1, 3.5]\,K. This comparison
uses the same averaging/reference periods but compares different inference
systems; it is not a numerical reproduction of Supplementary Figure S3.

\begin{figure}[htbp]
\centering\includegraphics[width=\textwidth]{climate_alternative.pdf}
\caption{Temperature-only calibration and independent observational checks.
The upper-right forecast includes residual and measurement uncertainty. The
lower-left scenario bands contain parameter uncertainty only. Ocean heat is
compared at the actual observational depth, with a training-period datum offset.
The temperature and heat holdout discrepancies prevent treating this converged
alternative posterior as validation of the target model.}
\end{figure}

\subsection{Joint temperature--heat calibration}

The next experiment adds independent NOAA 0--700\,m heat observations over
1955--2005, with a datum offset and separate AR1 residual parameters. The
first joint run failed convergence ($\widehat R=1.0218$ and a heat-correlation
ESS of about 285). Its diagnostics remain in the repository. A revised
sampler uses logarithmic innovation scales, logit correlations and a recent
warmup covariance window; Jacobians preserve the same physical priors.
Four chains of 8000 warmup and 50,000 production iterations retain 20,000
states. All acceptance checks pass: the maximum $\widehat R$ is 1.00582 and
the minimum bulk/tail ESS is 589.

An exact modal solver reduces computation while retaining all finite-volume
degrees of freedom. A diagonal similarity transforms the climate operator
into a symmetric tridiagonal matrix. Each mode is integrated analytically
under linearly interpolated annual forcing. Eighteen checks agree with an
independent augmented matrix exponential and establish second-order
Crank--Nicolson convergence. Four substeps per year differ from the exact
solution by at most 0.00133\,K in three historical comparisons; the exact
solver is used consistently for the new joint fit and projections.

The joint fit reduces held-out temperature RMSE to 0.0870\,K and bias to
$+0.0402$\,K. Heat predictive RMSE is $0.879\times10^{22}$\,J. All 16 held-out
years lie in both nominal 90\% predictive intervals. These short, correlated
holdouts provide limited validation; coverage is not a proof of correct
uncertainty. ECS has a median of 1.485\,K and a parameter 90\% interval of
1.093--1.958\,K. Future warming is substantially below the AR6 assessment:
SSP2-4.5 warming over 2081--2100 is 1.973 [1.518,2.483]\,K, compared with
2.7 [2.1,3.5]\,K; SSP5-8.5 is 3.319 [2.569,4.160]\,K, compared with
4.4 [3.3,5.7]\,K. This discrepancy remains unresolved. Improved historical
prediction does not justify replacing the target's joint climate--carbon--ice
calibration with this alternative.

\begin{figure}[htbp]
\centering\includegraphics[width=\textwidth]{climate_joint_alternative.pdf}
\caption{Alternative joint temperature and upper-ocean-heat calibration.
Observations after 2005 are withheld. Future bands propagate parameter
uncertainty under prescribed SSP forcing. The improved holdout performance
and the low late-century warming relative to AR6 are both retained.}
\end{figure}

\subsection{Antarctic events and observational thermal expansion}

DAIS geometry, accumulation and marine discharge are now implemented. The
discharge water depth follows sea level minus bed elevation, correcting the
printed sign from geometry; coefficient units follow dimensional consistency.
The equilibrium volume of $2.478479\times10^{16}$\,m$^3$ differs from Shaffer's
rounded $2.478\times10^{16}$\,m$^3$ by 0.0193\%. This is a reproduced
methodological benchmark, not a reproduced Wong figure.

The printed runoff gate activates at $h_0+cT_A=0$. Over marine bed this can
produce a discontinuous accumulation rate. Splitting time steps at that event
and evaluating endpoint stages on the appropriate side removes the earlier
irregular refinement. In the 6\,K, century-long warming ramp, the event occurs
at year 41.92982456. Two-year and one-year steps differ in year-200 radius by
$1.57\times10^{-7}$\,m; finer steps reach floating-point limits. The dashed
curves in Figure\ref{fig:recovered} instead activate runoff when the runoff
line exceeds local edge height, and remain explicit alternative physics.
These tests express sea level by mass conservation using ice/water densities
and ocean area, whereas B17 scales its initial Antarctic volume to 57\,m.
The two observable conventions are not silently interchanged.

\begin{figure}[htbp]
\centering\includegraphics[width=.98\textwidth]{recovered_mechanisms.pdf}
\caption{Independent mechanism tests. Left: prescribed Antarctic warming,
with fixed ocean temperature and sea level, isolates the runoff convention.
Right: climate grid refinement under 1\% annual CO$_2$ growth. Neither panel
is a global calibrated projection.}\label{fig:recovered}
\end{figure}

An independent coefficient check uses NOAA heat and thermosteric sea level
from the same depth range. OLS over 1955--2005 at 0--700\,m gives
$1.392\times10^{-25}$\,m\,J$^{-1}$; fitting 2005--2015 at 0--2000\,m gives
$1.107\times10^{-25}$\,m\,J$^{-1}$. Held-out thermosteric RMSEs through 2021
are 0.414 and 0.757\,mm, respectively. First-difference slopes provide a check
against merely fitting a common trend. Because heat and sea level share
hydrographic observations, these are descriptive constraints rather than
independent-error coefficient posteriors. They are not substituted silently
for the unavailable V19 prior or a full-depth expansion coefficient.

\begin{figure}[htbp]
\centering\includegraphics[width=.98\textwidth]{thermal_observations.pdf}
\caption{Independent NOAA heat and thermosteric observations with fitted
depth-specific linear relations. Later years through 2021 were withheld.
This coefficient check is an explicit alternative to the target calibration.}
\end{figure}

\subsection{Coastal cash flows and input quality}

The annual coastal experiment retains the documented ten-choice menu and
first-period decision rule. Sixty checks establish nonnegative cash-flow
components, first-period policy selection, no repeated relocation at fixed
sea level and telescoping incremental dike construction costs. Upgrades at
2010, 2051 and 2101, annual payment timing and a paired no-SLR counterfactual
are explicit alternatives to the target's decadal accounting and initial-step
subtraction. Forty-eight synthetic scenarios compare income/exposure, sea-level
rise, retreat land accounting and linear versus quadratic maintenance.
All exposure values are labeled fixtures; no global cost estimate is inferred.
Doubling flood quadrature from 512 to 1024 panels changes the verification-case
NPV by about 0.00024\,USD and leaves its policy unchanged.

\begin{figure}[htbp]
\centering\includegraphics[width=.98\textwidth]{coastal_alternative.pdf}
\caption{Synthetic coastal alternatives. Left: paired-counterfactual incremental
costs under two explicit land-accounting rules. Income and exposure both vary
across the three fixtures, so this is not an isolated income effect.
Right: annual components for one selected retreat policy, including discrete
upgrades. These outputs do not estimate a named coast or a global cost.}
\end{figure}

The independent GTSR shapefile has 12,148 coastal feature rows but only 7868
distinct station labels. The initial interpretation of these labels as unique
DIVA identifiers was corrected by this audit. A second public release provides
16,611 station-wise Gumbel fits. Geographic mapping leaves 133 nearest matches
with invalid parameters; these remain missing, not filled from distant stations.
The two input products also disagree in their 100-year levels. These findings
prevent a silent merge or assumption of exact input equivalence. The full
target elevation-area, exposure and DINAS tide records still require recovery.
The next section uses a separately sourced geographic input product for an
explicit alternative, with different segmentation and monetary conventions.

Every alternative here has its own input, script, output and figure record.
No alternative curve is entered into the target inventory as a recovered
published numerical result. Ongoing calibration and data-recovery tasks remain
active rather than being reclassified as irreducible merely because the first
implementation lacked them.
